\section{Discussion and limitations}
\label{sec:discussion}

\paragraph{Whole-document replay}
The manifest detects the deletion, reordering, duplication or individual replay of an encrypted field within a resource. It does not detect the replacement of an entire resource, manifest included, by an older version of itself: the old resource is internally consistent and was genuinely produced for this patient, tenant and type. Preventing that requires a monotonic version, a sequence number or a commit timestamp, to be part of the context, and the version is application state the library cannot see. An application that needs it puts the version in the \code{purpose} field or extends the context, and then must supply the current version at every read, which is the same discipline the rest of the design already demands.

\paragraph{The context must be authentic}
Everything in Section~\ref{sec:design} rests on the caller supplying the context from authenticated facts. The API takes the context as a required argument to every operation and stores none of it; the header carries no context field; the FHIR layer may derive the patient from the resource at encryption time but requires it from the caller at decryption. These choices make the mistake visible but cannot make it impossible. An application that stores the context beside the ciphertext and reads both back has rebuilt the shared-key baseline, because the adversary who can move the ciphertext can move the stored context with it. This is the single rule a deployment has to get right, and it is stated in the library's documentation before the first line of usage.

\paragraph{Metadata}
Ciphertext lengths, the existence of a record, its write time, and which fields of a resource are encrypted are all visible to storage. Length-hiding padding to a bucket size is compatible with the construction and was not a goal. An adversary who can count encrypted notes learns that a patient has that many notes.

\paragraph{Out of scope}
Compromise of the application host while master keys are in memory defeats any scheme at the storage layer; the defence is host hardening and hardware key storage, which the key-wrapping interface accommodates but this paper does not evaluate. Timing side channels beyond constant-time comparison of the commitment are delegated to the underlying libraries. Availability is not addressed: an adversary who deletes storage denies service, and the construction guarantees only that whatever storage returns is either correct or rejected.

\paragraph{Relation to the AWS Database Encryption SDK}
The deployed construction closest to CPDPM is the AWS Database Encryption SDK (Section~\ref{sec:practice}), and the agreements are substantive: client-side encryption of individual fields, a fresh data key per encryption, the record's identity in an encryption context the ciphertext is bound to, a signature over the record that detects field addition, deletion and value swapping, key commitment, and a hierarchical keyring with versioned branch keys and one branch key per tenant~\cite{aws-dbesdk-concepts,aws-dbesdk-hierarchical}. The differences follow from the data model and the trust placed in the host. The SDK binds to a DynamoDB item's primary key, so a value moved between two attributes of the same item is caught by the item signature rather than by the binding; CPDPM binds to tenant, patient, resource type and field path, with an encoding specified so that contexts cannot collide, and the binding alone rejects the move. The SDK is specific to DynamoDB and to its attribute model; CPDPM stores a self-describing string in any column, document or object store, and its manifest is shaped by FHIR rather than by a table. The SDK protects each data key by envelope encryption, storing it wrapped in a material-description field of the record, at the cost of a wrapping operation and some bytes per record, and its hierarchical keyring caches plaintext branch keys on the host; CPDPM derives per-context keys on demand from a master key through HKDF and stores an eleven-byte header, with the master key resident on the host throughout. The SDK's optional ECDSA signature separates writers from readers, which CPDPM, symmetric throughout, does not. A deployment on DynamoDB without FHIR structure should use the SDK. The contribution here is the same discipline made database-agnostic and FHIR-shaped, with its binding and exposure claims measured against the baseline that practice already reaches.

\paragraph{Review status}
The security argument of Section~\ref{sec:security} is a sketch reducing to standard notions, not a formal proof. The implementation was reviewed by its author against the threat model, with property-based fuzzing of every parser, and the findings and fixes of that review are recorded in the artifact. Two of them were substantive: field deletion and per-field replay within a FHIR resource were not detected until the manifest was added. No third party has examined the code, and the construction should not protect real patient data until one has.

\paragraph{Unicode}
Context strings are normalised to NFC before encoding. Injectivity of the encoding does not depend on normalisation: it is a property of the length-prefixed byte string, holds for any inputs, and is property-tested. What normalisation affects is whether two interpreters canonicalise the \emph{same} identifier to the same bytes, and a disagreement there causes a false rejection, never a false acceptance. The Unicode normalisation stability policy guarantees that a string of assigned characters normalises identically under every later version of the standard~\cite{unicode-stability}, so the only source of disagreement is a code point unassigned on one interpreter and assigned on another. The library therefore rejects unassigned code points at context construction, alongside control characters and surrogates, which makes validity depend on the interpreter's Unicode tables only for characters that no identifier should contain.

\paragraph{Compliance}
The introduction motivates the work from regulation, so it should say what a compliance review would find. The AES-256-GCM suite, HKDF-SHA256 and HMAC-SHA256 are FIPS-approved algorithms and are called through OpenSSL, for which a FIPS 140-3 validated provider exists~\cite{openssl-fips}; the pre-built wheels of the \code{cryptography} package link OpenSSL statically and make no validation claim~\cite{cryptography-install}, so a deployment that needs a validated module must build the package against a system OpenSSL configured with the FIPS provider. XChaCha20-Poly1305 and Argon2id are not FIPS-approved, which is why the AES suite is the default and the PBKDF2 path exists. HKDF with an empty salt is the construction RFC~5869 specifies for uniformly random input keying material~\cite{rfc5869}. None of this makes a deployment compliant; it tells the officer which boxes the construction can tick and which the deployment has to.

\paragraph{Performance}
The fixed per-operation cost of two HKDF derivations and a commitment is about \CpdpmFixedUs\,$\mu$s on the arm64 machine and \XCpdpmFixedUs\,$\mu$s on x86-64, which matters only for workloads encrypting many small fields at high rate, and there per-field binding is the point. The XChaCha20-Poly1305 suite is slow in the current build for reasons of library and hardware support rather than design, and a ChaCha20-Poly1305 suite over the same OpenSSL bindings as AES-GCM is the natural addition.

\paragraph{Reading the hybrid-encryption literature}
Section~\ref{sec:pretransform} suggests three questions a reviewer can put to any scheme that adds a layer in front of a standard cipher. Does the layer's key travel with the ciphertext? Is there a reduction, or only a statistic of the output? And what can an adversary do to the scheme that the same adversary cannot do to the cipher alone? For the twelve schemes in Table~\ref{tab:genre}, the evidence reported is statistical, none states a reduction, and where key transport is described the key is carried with the data or derived from it; the third question is the one this paper answers.
